\section{Conclusion}
\label{sec:conclusion}

We introduced \ours{}, a minimal modification to standard deep ensembling:
train the same architecture with the same cosine schedule but two different
$T_\mathrm{max}$ endpoints, and ensemble the resulting checkpoints. The two
trajectories reach identical mean accuracy yet differ on 13\% of test samples,
producing a $+1.5$~pp gain beyond standard same-length deep ensembles.

The trick is validated across 6 datasets spanning two modalities and three
architectures. It achieves a new state of the art on FACED 9-class EEG
emotion classification ($\facedacc$ BAcc, $+8.2\%$ over the prior
EMOD AAAI 2026 baseline of $0.642$). The gain scales inversely with
individual accuracy: capacity-bound students (FACED 9-class, SEED-V
5-class) see the largest gains ($+3$--$4$~pp), easy tasks (MNIST) see
minimal gains ($+0.14$~pp).

The theoretical mechanism is transparent: the extra epochs in the
longer trajectory constitute a ``low-LR exploration'' phase that drifts
weights into a neighboring local minimum with identical loss but different
decision boundaries. The 100ep and 150ep consensus predictions differ on
13\% of test samples, and on those samples the mixed ensemble is 10~pp more
accurate than either individual consensus---directly predicting the overall
$+1.5$~pp gain.

\ours{} is simple enough to deploy in any existing cosine-schedule
training recipe with one line of code (double one integer) and provides
consistent calibration-preserving gains across tasks. We hope it will
become a standard entry in the practitioner's toolkit of ``free'' ensemble
tricks.
